\chapter{Introduction}

Tumour heterogeneity is a major driver of drug resistance and treatment failure. Rare subpopulations of
tumour cells may underlie resistance \cite{scp542}, the few cells that survive treatment can ultimately
cause relapse \cite{scdeal}, and in one clinical trial the predicted response of the most resistant clone
separated responders from non-responders better than an average over clones \cite{perception}. Drug
response, however, is measured on bulk populations of cells, and no large-scale screen measures it per
cell (Sections~\ref{sec:bg-screens} and~\ref{sec:bg-singlecell}). Single-cell RNA sequencing (scRNA-seq)
resolves the subpopulations that bulk measurements average over \cite{luecken2019,heumos2023}, so a model that
predicts response from single cells could report a distribution of sensitivities within a sample rather
than one average \cite{perception,scdeal}.

This report trains such a model on the labels that exist. The response labels of the Cancer Therapeutics
Response Portal v2 (CTRPv2) \cite{ctrpv2,ctrp_rees} are measured on bulk cell populations. They are mapped onto the individual
cells of the corresponding cell lines in the SCP542 pan-cancer atlas \cite{scp542}, so that one value per
(cell line, compound) is broadcast onto a few hundred heterogeneous cells, and the setup is therefore
weakly supervised.

Cell lines do not fully recapitulate tumours, but they provide the volume of
continuous labels needed to train a representation, and they are heterogeneous in expression
(Section~\ref{sec:bg-heterogeneity}) \cite{scp542}.

The cost is a mismatch of resolution. One bulk value per line adds label noise to every
cell, and because all cells of a line share it, they are pseudo-replicates, so the effective sample size
is the number of cell lines rather than the number of cells \cite{lazic2010,dreval}.

The cells are represented in two ways, by principal component analysis (PCA) fitted to the atlas
\cite{jolliffe2016,luecken2019} and by the frozen embedding of scGPT, a transformer pretrained on about
33 million human cells \cite{scgpt} (Section~\ref{sec:bg-foundation}). The hypothesis behind the
comparison was that scGPT provides a denoised biological prior which aligns cellular states across
tissues, so that a regressor is encouraged to learn transcriptomic drivers of sensitivity rather than to
memorise the tissue or line of origin. Single-cell counts, by contrast, are sparse, with many zeros due to
sampling \cite{luecken2019}, and the premise was that regressing them directly onto bulk labels would be
unstable. In the latent space of the atlas, PCA separates individual cell lines into compact islands,
whereas scGPT places the same cells on one continuous manifold (Figure~\ref{fig:umap}), a contrast at the
granularity at which the response label is defined.

This design addresses three
questions: whether drug response can be predicted from single cells with bulk labels at all, whether
scGPT embeddings are a viable representation for this task against PCA, and whether a model trained on
single cells learns cellular heterogeneity implicitly (Section~\ref{sec:bg-objective}).


The scope of this report is a single database, a single continuous response score and a supervised
regressor trained on a panel of \NPanel{} compounds. Because DrEval shows that published drug-response
models often score well through compound and cell-line mean effects alone \cite{dreval}, the panel scores
are read against a ridge regression on cell-line means \cite{hoerl1970}, against DrEval's naive baselines
and its per-compound random forest, against the same model trained on all \NDrugsModelled{} compounds, and
with general cell-line sensitivity removed. The comparison of the two representations is tested against
five alternative explanations, and the heterogeneity analysis against a positive control and technical
confounders. Cross-database integration, toxicity and clinical fine-tuning are left to later work.
